% Agradecimientos: va después de la advertencia y antes del índice, en una hoja sin encabezado ni número
\thispagestyle{empty}

\vspace*{1cm}
{\centering \Large \textbf{Agradecimientos}\par}
\vspace{0.8cm}

% Familia
Quiero expresar mi más profundo agradecimiento a mi madre, Patricia Lopez Zamora; a mi padre, Juan Jesús Angeles Gonzales, y a mi hermana, Jessica Vianey Angeles Lopez, por su apoyo incondicional a lo largo de mi carrera. Gracias por motivarme a seguir adelante, por ayudarme a cumplir mis sueños, por estar siempre a mi lado y por nunca rendirse conmigo, incluso en los momentos más difíciles. A ustedes les debo cada uno de mis logros y les estaré eternamente agradecido.

% Director
Agradezco a mi director, el Dr. Genaro Juárez Martínez, por su invaluable guía durante el desarrollo de este proyecto y a lo largo de mi carrera. Gracias por introducirme en el mundo de la vida artificial y el estudio de los sistemas complejos, por motivarme a participar en proyectos en el extranjero y por todas sus enseñanzas. Su orientación fue fundamental para dar rigor a este trabajo.

% Mentores
Agradezco también a Willem Nielsen, mi mentor en la \textit{Wolfram Summer School}; a Stephen Wolfram, por sus charlas sobre autómatas celulares y sus valiosos consejos, y a Sukanta Das, por su colaboración en la redacción de un artículo durante mi estancia en la India.

% Compañeros
Gracias a todos los amigos que encontré en este camino: a mis compañeros de la ESCOM, de la \textit{Wolfram Summer School}, de Infosys y de \textit{Cellular Automata India}. Cada momento, cada plática y cada experiencia compartida hicieron de este viaje algo único.

% Comunidades de código abierto
Este proyecto no habría sido posible sin los recursos de código abierto y contenido libre que distintas comunidades y grupos de investigación desarrollan y mantienen de forma continua. Agradezco su trabajo y su compromiso con el conocimiento compartido.
\clearpage
